%% ma: L10-B14-0025
%% lop: 10
%% chuong: 5
%% bai: 14
%% dang: 10.14.5
%% loai: DS
%% muc: [NB, TH, TH, VD]
%% dap_an: "ĐSĐĐ"
%% nguon: "chuyen-de-cac-so-dac-trung-cua-mau-so-lieu-khong-ghep-nhom-toan-10-chuong-trinh-moi.pdf (GV. Phan Nhật Linh) (P04), Bài 03 Dạng 1 Phần II Câu 2"
%% kien_thuc: "Lớp 10 Bài 13 (tứ phân vị), Bài 14 (khoảng tứ phân vị, giá trị bất thường)"
%% trang_thai: chinh-thuc
%% ghi_chu: "Đề gốc là giá của 8 loại rượu ngoại, đã đổi thành 'loại hàng' (số liệu giữ nguyên). Tự duyệt theo yêu cầu giáo viên 22/10/2026."
\begin{ex}
Mẫu số liệu sau là giá tiền (đơn vị: triệu đồng) của $8$ loại hàng được nhập về tại một cửa hàng:
\begin{center}
$1{,}2$\quad $1{,}35$\quad $1{,}42$\quad $1{,}53$\quad $1{,}8$\quad $1{,}84$\quad $1{,}96$\quad $2{,}4$
\end{center}
Xét tính đúng sai của các khẳng định sau:
\choiceTF{\True Khoảng biến thiên của mẫu là $R=1{,}2$}{Khoảng tứ phân vị là $\Delta_Q=0{,}215$}{\True $Q_1-1{,}5\Delta_Q=0{,}6125$}{\True Mẫu số liệu không có giá trị bất thường}
\loigiai{a) Đúng: $R=2{,}4-1{,}2=1{,}2$.

b) Sai: $Q_1=\dfrac{1{,}35+1{,}42}2=1{,}385$; $Q_3=\dfrac{1{,}84+1{,}96}2=1{,}9$; $\Delta_Q=0{,}515$.

c) Đúng: $1{,}385-1{,}5\cdot0{,}515=0{,}6125$.

d) Đúng: $Q_3+1{,}5\Delta_Q=2{,}6725$; mọi giá trị đều thuộc $[0{,}6125;2{,}6725]$.}
\end{ex}
